\section{Design Philosophy and Limitations}

\subsection{Design philosophy}

The design favors clarity and incremental learning. Each advanced feature is introduced in a separate variant rather than hidden inside one large, configurable CPU.

This is a good approach for documentation because it lets readers compare:

\begin{itemize}
\item single-cycle vs multi-cycle timing;
\item non-pipelined vs pipelined hazards;
\item predict-not-taken vs dynamic branch prediction;
\item physical addressing vs MMU translation;
\item simple trap handling vs precise pipelined traps.
\end{itemize}
\subsection{Known simplifications}

Important simplifications include:

\begin{itemize}
\item combinational divider in \texttt{alu.v};
\item educational combinational MMU in \texttt{mmu.v};
\item no TLB;
\item limited privilege model;
\item incomplete production-grade CSR set;
\item atomics not uniformly integrated into every variant;
\item instruction-side MMU behavior is not as complete as a production Linux-capable core;
\item branch predictor is small and direct-mapped by PC index;
\item no cache hierarchy.
\end{itemize}
These are reasonable for a from-scratch learning CPU. They also clearly identify future improvement points.

\bigskip
